% !TEX root = ../main.tex
%%%%%%%%%%%%%%%%%%%%%%%%%%%%%%%%%%%%%%%%%%%%%%%%%%%%%%%%%%%%%%%%%%%%%%%
% Calibration and anchoring
%%%%%%%%%%%%%%%%%%%%%%%%%%%%%%%%%%%%%%%%%%%%%%%%%%%%%%%%%%%%%%%%%%%%%%%

\section{Anchoring and calibration}
\label{sec:calibration}

A model of this kind can be calibrated in two quite different senses, and
conflating them is how design targets come to be reported as results. We
therefore assign every parameter to one of three \emph{anchoring tiers}, and
state the tier wherever the parameter is used.

\begin{enumerate}[label=\textbf{(\arabic*)}]
  \item \textbf{Anchored.} There exists an empirical quantity of which the
  parameter is an estimate, a citable source for it, and a comparator defined
  on the same perimeter as the model. $\sigma$, $\pi_0$, $\eta$, $\lambda$ and
  the replacement rate are in this tier.
  \item \textbf{Declared convention.} A value is needed, the literature offers
  a range rather than a point, and the choice within that range is ours. These
  are not estimates and are never described as such. $\delta$ is the
  consequential case.
  \item \textbf{Non-anchorable by decision.} No empirical quantity corresponds
  to the parameter at all, so searching for a citation would be cosmetic.
  $c_0$ is the case, and it is handled by reporting both regimes throughout
  rather than by choosing.
\end{enumerate}

\Cref{tab:params} lists every parameter with its value and its tier. The tiers
are not decoration: they record which numbers carry empirical content and which
are modelling decisions that must not be presented as estimates.

\begin{table}[htbp]
\centering
\small
\caption{Parameters, values and anchoring tiers.}
\label{tab:params}
\begin{threeparttable}
\begin{tabularx}{\textwidth}{llX}
\toprule
Parameter & Value & Anchoring \\
\midrule
\multicolumn{3}{l}{\textit{Anchored to a citable source}}\\
\addlinespace[2pt]
$\sigma$ (elasticity of substitution) & swept & $0.40$--$0.60$
  \citep{Chirinko2008}; $0.45$--$0.87$ and Cobb-Douglas rejected
  \citep{Knoblach2020}. \textbf{Swept, never chosen.}\\
$\pi_0$ (base-point capital share) & $1/3$ & standard growth accounting
  \citep{Gollin2002,Solow1957}\\
$\eta$ (wage-curve elasticity) & swept, $0$--$0.15$ & $\approx 0.10$
  \citep{BlanchflowerOswald1994}; $\approx 0.07$ \citep{NijkampPoot2005}\\
$\lambda$ (wealth effect) & $0.05$ & $\approx 5$ cents, 16-country mean
  \citep{Slacalek2009}; cf.\ $3$--$4$ cents for housing wealth
  \citep{CaseQuigleyShiller2005}\\
$\rho$ (retention ratio) & $0.40$ & anchored \emph{through} $I/Y$, not through
  payout ratios; see \cref{sec:limits}\\
$rr$ (benefit replacement rate) & swept, $0$--$0.75$ & OECD net replacement
  rates $\approx 58\%$ initial \citep{OECD2024}\\
\addlinespace[4pt]
\multicolumn{3}{l}{\textit{Declared conventions --- not estimates}}\\
\addlinespace[2pt]
$\delta$ (depreciation) & $0.05$ & BEA-implied figure for this perimeter is
  $\approx 0.090$; structures alone $2$--$3\%$. A defensible midpoint,
  \textbf{not a measurement}. See \cref{sec:delta}.\\
$U_{\min}$, $w_{\min}$ & $0.01$, $0.45$ & guards; $w_{\min}$ never stably
  binds anywhere tested\\
$\max\tau$ (tax cap) & $0.6$ & guardrail; realised $\tau$ is a measured
  outcome\\
$\underline{I}$ (investment floor) & $0.1$ & numerical guardrail\\
$\bar{u}$ (target utilisation) & $0.90$ & above the observed US average
  ($\approx 0.80$)\\
\addlinespace[4pt]
\multicolumn{3}{l}{\textit{Modelling choices --- measured once, then frozen}}\\
\addlinespace[2pt]
$(K_0,L_0)$ & $41.87$, $7.395$ & CES normalisation anchor
  \citep{KlumpDeLaGrandville2000}; caveat per \citet{Temple2012}\\
$U_{\mathrm{REF}}$ & $0.26047$ & wage-curve normalisation point.
  \textbf{Not an estimate of the NAIRU.}\\
$\bar{w}$ & $0.9$ & demoted to a normalisation point once $\eta$ exists\\
$K_{\text{init}}$ & $40.0$ & \textbf{selects the basin} (multiple equilibria);
  held fixed in every sweep\\
$\beta$ (accelerator) & $0.5$ & no canonical value exists; treated by
  sensitivity analysis\\
\addlinespace[4pt]
\multicolumn{3}{l}{\textit{Declared non-anchorable}}\\
\addlinespace[2pt]
$c_0$ (autonomous consumption) & $\{1.0, 2.0\}$ & model units have no exchange
  rate with data (numéraire $=1$). Its original justification was
  \emph{falsified}; both regimes are reported everywhere.\\
\bottomrule
\end{tabularx}
\begin{tablenotes}[flushleft]\footnotesize
\item \textit{Note.} One model period $=$ one year. This is a declared
interpretive convention, chosen as the only one consistent with the already
anchored parameters: $\delta$ is an annual rate and the Blanchflower--Oswald
elasticity is estimated on annual data. It follows that the 2{,}000 steps of
each run are a \emph{convergence device, not 2{,}000 years of simulated
history} --- every result in this paper is comparative statics across steady
states, never a time series to be overlaid on data.
\end{tablenotes}
\end{threeparttable}
\end{table}

\subsection{The capital block, anchored flows first}

The capital block is anchored \emph{flows first}: $\rho$ is anchored through
the investment rate it produces rather than through corporate payout ratios,
because $\rho$ is what fixes $I/Y$ here. The comparator is declared and held on
both sides of the ratio --- private \emph{non-residential} fixed capital ---
since mixing perimeters is exactly how an earlier version of these notes
derived a $\delta$ that appeared anchored and was not.

Against the BEA series for that perimeter, $I/Y$ is $0.138$--$0.141$
\citep{FRED_PNFI}. The model \emph{measured} at $\rho = 0.40$ gives $0.158$ in
the anchor scenario and $0.182$ in the headline scenario --- above the anchor by
two and four points. We quote the measured figure and never an analytical one:
the relation $I/Y = \rho\alpha$ belongs to an earlier Cobb-Douglas core without
a labour market and no longer holds once profit is a residual
(\cref{app:retractions}).

\subsection{The anchored retention rate}
\label{sec:anchoredrho}

Several statements in this paper turn on where the \emph{empirically anchored}
$\rho$ sits, so it is derived here rather than assumed. The procedure is the
one just stated, read backwards: $\rho$ is anchored through $I/Y$, so the
anchored $\rho$ is the value at which measured $I/Y$ enters the BEA band, by
linear interpolation between the bracketing nodes of the committed sweep
(\cref{tab:anchoredrho}).

\begin{table}[htbp]
\centering
\small
\caption{The anchored retention rate, derived from measured $I/Y$ against the
BEA band $0.138$--$0.141$.}
\label{tab:anchoredrho}
\begin{tabular}{llll}
\toprule
Scenario & $I/Y$ at $\rho = 0.35$ & Anchored $\rho$ & Status\\
\midrule
Anchor ($c_0 = 2.0$, $\sigma = 1$) & 0.135 & $0.357$--$0.363$ & interpolated
in support\\
Headline ($c_0 = 1.0$, $\sigma = 0.5$) & \textbf{0.161} & $< 0.35$ & below the
swept support\\
\bottomrule
\end{tabular}
\end{table}

The two scenarios do not give the same number, and the difference cuts in the
paper's favour rather than against it. In the anchor scenario the band is
reached inside the support, at $\rho \approx 0.36$. In the headline scenario
measured $I/Y$ is already $0.161$ at $\rho = 0.35$, the \emph{lowest} value
swept --- so the anchored $\rho$ lies below the support entirely. We decline to
extrapolate it: a value obtained by extending the fit past the sweep would be a
number about a region never visited.

What both scenarios support is an upper bound, and the bound is what the
argument uses: \textbf{the anchored $\rho$ is at most $0.363$}. Wherever this
paper says the anchored retention rate lies to the left of a turning point, the
claim is that the turn lies above $0.363$, which is checkable against any table
of $\rhostar$ and does not depend on resolving the anchored value exactly.

\begin{declbox}
An earlier version of this project carried the anchored range as
$\rho \in [0.35, 0.45]$ --- the lower end derived as above, the upper end
\emph{not derived at all}. The reported figures below use the derived bound.
The committed sensitivity artifact still carries a side column computed under
the wider band and therefore labels two low-$\sigma$ cells as straddling the
turn; \cref{tab:turning} reports the corrected side, and the discrepancy is
recorded rather than regenerated, since the artifact is evidence of what was
run.
\end{declbox}

$\delta$ is the clearest case of tier~(2). The depreciation rate implied by the
BEA data for this perimeter is $\approx 0.090$ \citep{FRED_Depreciation},
inflated by intellectual
property products; structures alone sit at $2$--$3\%$, so the defensible
empirical band spans roughly $0.02$--$0.09$ and $\delta = 0.05$ is a midpoint
chosen within it. It is not a measurement, and it must not be recalibrated
``toward the centre of the literature'' late in a project: doing so would move
every canonical number without buying any anchoring.

That decision has a structural consequence which \cref{sec:delta} makes
quantitative. The steady-state closure is $I/K = \delta + g$. This model has
$g = 0$, so $K/Y$ is not an anchor at all but a \emph{mechanical outcome}:
$K/Y = (I/Y)/\delta$, and measured $I/K$ is $0.0500$, exactly $\delta$. Model
$K/Y$ is therefore $3.17$--$3.64$ against a business-sector $1.23$
\citep{FRED_NetStock}, and no choice of $\rho$ reconciles the two. \emph{The model cannot match business
$I/Y$ and $K/Y$ simultaneously without trend growth} --- a declared structural
limitation, not a calibration failure to be tuned away.

\subsection{What is deliberately not anchored}

$c_0$ is declared \emph{non-anchorable by decision}: it is autonomous
consumption in model units, and since the numéraire is~1 by construction there
exists no empirical quantity of which it is an estimate. Searching for a
citation would be cosmetic anchoring. Its sensitivity is instead handled by
reporting both regimes $c_0 \in \{1.0, 2.0\}$ throughout, and by the global
analysis of \cref{sec:sa}. This does not close the model's most visible
quantitative gap --- unemployment rates far above anything observed --- which is
a consequence of design decisions listed in \cref{sec:limits}, not of an
unanchored parameter.
